\chapter{Monotone solutions and minimizing properties}\label{sec:phase-minimizers}
Stability is a second-variation condition. A minimizing property compares finite changes and needs an additional argument. Here $\varepsilon=1$, $W(s)=(1-s^2)^2/4$, and $n\geq2$. We develop the sliding proof in \cite{ChAC}, Section 5.1, pp.~17--18.

\begin{definition}[Minimization on bounded balls]\label{def:phase-ball-minimizer}
A smooth entire function $u$ minimizes the Allen--Cahn energy on bounded balls if, for every Euclidean ball $B$ and every $v\in H^1(B)\cap L^4(B)$ with $v-u\in H^1_0(B)$,
\[
 E_1(u;B)\leq E_1(v;B).
\]
In particular it minimizes against every perturbation supported in a compact subset of space. This comparison is stronger than the nonnegativity of its quadratic second variation. It is the minimizing property used in \cite{ChAC}, Proposition 5.8.
\end{definition}

\begin{lemma}[Strong contact for ordered solutions]\label{lem:phase-ordered-contact}
Let $u,v$ be smooth solutions of $\Delta u=W'(u)$ and $\Delta v=W'(v)$ on a connected open set. If $u\leq v$ and equality holds at one interior point, then $u=v$ throughout the set.
\end{lemma}
\begin{proof}
The function $w=v-u\geq0$ solves $\Delta w=c(x)w$, where
$c(x)=\int_0^1W''(u+t(v-u))dt$ is locally bounded. On a relatively compact ball choose $C\geq\max\{0,\sup c\}$. Then $z=-w\leq0$ satisfies $(\Delta-C)z=(C-c)w\geq0$. At a zero of $w$, $z$ attains the nonnegative maximum zero. The strong maximum principle, whose zero-order coefficient is now nonpositive, makes $z=0$ on that ball. Thus the zero set of $w$ is open as well as closed; connectedness finishes the proof. This proves \cite{ChAC}, Lemma 5.7, without assuming an incorrect sign for $W''$.
\end{proof}

\begin{theorem}[A bounded Dirichlet minimizer on a ball]\label{thm:phase-dirichlet-minimizer}
Let $B$ be a bounded ball and $g\in C^\infty(\overline B)$ with $-1<g<1$ on $\partial B$. There exists a minimizer of $E_1(v;B)$ over $v\in H^1(B)\cap L^4(B)$ with $v-g\in H^1_0(B)$. One can choose it smooth in $B$, continuous on $\overline B$, with $-1<v<1$ everywhere on $\overline B$. No uniqueness is asserted.
\end{theorem}
\begin{proof}
The class is nonempty, and its infimum is finite and nonnegative. A minimizing sequence has bounded gradient energy and, by quartic coercivity, bounded $L^4$ norm. It is therefore bounded in $H^1(B)$. Weak compactness and Rellich compactness on the ball give a subsequence converging weakly in $H^1$, strongly in $L^2$, and almost everywhere. The affine zero-trace class is weakly closed. Fatou gives lower semicontinuity of the nonnegative potential and finite $L^4$ norm of the limit; weak lower semicontinuity handles the gradient. Thus a minimizer exists.

Truncation to $[-1,1]$ preserves the boundary data and cannot increase energy, so choose a bounded minimizer. Compactly supported variations give $\Delta v=W'(v)$ weakly. Interior elliptic regularity makes it smooth. Boundary regularity for the Poisson equation on the smooth ball with bounded forcing and smooth boundary data gives continuity on $\overline B$ and the prescribed boundary values. These are external elliptic and compactness inputs, not consequences of weak convergence alone. Comparison with the constant solutions $\pm1$ via Lemma \ref{lem:phase-ordered-contact} gives strict inequalities in the interior; a constant well would contradict the boundary data. Together with the strict boundary inequalities this proves the assertion. This fills in the existence step of \cite{ChAC}, Exercise 5.2(a).
\end{proof}

\begin{lemma}[Uniform limits of monotone translates on compact sets]\label{lem:phase-translate-limits}
Suppose $u$ is continuously differentiable on $\mathbb R^n$, $\partial_{x_n}u>0$, and
$u(x',t)\to\pm1$ as $t\to\pm\infty$ for every $x'$. Then $u(x+se_n)\to\pm1$ uniformly on every compact subset as $s\to\pm\infty$.
\end{lemma}
\begin{proof}
For integer $s\to+\infty$, these continuous functions increase pointwise on a compact set to the continuous constant one. Dini's theorem gives uniform convergence. Monotonicity bounds intermediate real parameters between consecutive integers. Apply the decreasing version at minus infinity. This justifies the compact uniform convergence used in \cite{ChAC}, Proposition 5.8.
\end{proof}

\begin{theorem}[Monotonicity with saturated ends gives minimization]\label{thm:phase-monotone-minimizing}
Let $u:\mathbb R^n\to(-1,1)$ be a smooth solution of $\Delta u=W'(u)$ with
\[
 \partial_{x_n}u>0,\qquad
 \lim_{t\to\pm\infty}u(x',t)=\pm1
 \quad\text{for every }x'\in\mathbb R^{n-1}.
\]
Then $u$ minimizes energy on bounded balls in the sense of Definition \ref{def:phase-ball-minimizer}.
\end{theorem}
\begin{proof}
Fix a ball $B$ and choose the minimizer $v$ with boundary data $u$ from Theorem \ref{thm:phase-dirichlet-minimizer}. Compactness of $\overline B$ and strict bounds give $\max v<1$ and $\min v>-1$. By Lemma \ref{lem:phase-translate-limits}, $u_s(x)=u(x+se_n)$ satisfies $u_s\geq v$ on $\overline B$ for sufficiently large positive $s$.

Let $s_*\geq0$ be the infimum of the parameters $s\geq0$ with this property. Monotonicity makes the admissible parameters an upper interval, and continuity gives $u_{s_*}\geq v$. If $s_*>0$, strict ordering on $\partial B$ follows from $u_{s_*}>u=v$ there. Unless equality occurs somewhere in $B$, compactness gives a positive gap on $\overline B$, permitting a slightly smaller shift, a contradiction. Interior contact forces $u_{s_*}=v$ on $B$ by Lemma \ref{lem:phase-ordered-contact}, again contradicting the strict boundary ordering. Hence $s_*=0$ and $u\geq v$.

Start instead with large negative translates, which lie below $v$, and move them upward toward shift zero. The same argument, with the ordering reversed, yields $u\leq v$. Thus $u=v$ on $B$, and it attains the Dirichlet minimum. Since the ball was arbitrary, the claimed property holds. This proves \cite{ChAC}, Proposition 5.8; the two limiting phases are a separate hypothesis from stability.
\end{proof}

\begin{proposition}[Area-order energy bound for bounded minimizers]\label{prop:phase-minimizer-growth}
If a smooth entire solution $u:\mathbb R^n\to[-1,1]$ minimizes energy on bounded balls, then
\[
 E_1(u;B_R(a))\leq C(n)R^{n-1}\quad(R>0)
\]
for every center $a$.
\end{proposition}
\begin{proof}
Lemma \ref{lem:phase-entire-estimates} gives a universal gradient bound. For $R\geq2$, choose a radial cutoff $\eta$ equal to zero on $B_{R-1}(a)$ and one near $\partial B_R(a)$, with $|D\eta|\leq C$. The competitor $v=1+\eta(u-1)$ has the same trace, takes values in $[-1,1]$, and has bounded gradient. It equals a zero-energy well in the inner ball, so its energy is at most a fixed constant times the volume of the width-one annulus, at most $C(n)R^{n-1}$. Minimality gives the same bound for $u$. For $R<2$, the bounded energy density gives $E_1(u;B_R(a))\leq CR^n\leq2CR^{n-1}$. This proves \cite{ChAC}, Exercise 5.4. In dimension three the bound meets the growth hypothesis of Theorem \ref{thm:phase-stable-growth-rigidity}, since minimization implies stability by second variation.
\end{proof}

\begin{proposition}[Limits at the ends of a monotone solution]\label{prop:phase-monotone-end-limits}
Suppose $u:\mathbb R^n\to(-1,1)$, $n\geq2$, is a smooth Allen--Cahn solution with $\partial_{x_n}u>0$. Then
\[
 u(x+te_n)\longrightarrow u_\pm(x')\quad\text{in }C^\infty_{\mathrm{loc}}(\mathbb R^n)
 \quad\text{as }t\to\pm\infty.
\]
The functions $u_\pm:\mathbb R^{n-1}\to[-1,1]$ solve the Allen--Cahn equation and are stable in dimension $n-1$. Their classification is a separate lower-dimensional problem.
\end{proposition}
\begin{proof}
Bounded monotonicity gives pointwise limits, and a fixed shift in $x_n$ does not change either limit, so they depend only on $x'$. Lemma \ref{lem:phase-entire-estimates} gives smooth local precompactness of the translates. Every subsequential limit must equal the same pointwise limit. The subsequence criterion therefore gives full smooth local convergence, and the equation passes to the limit.

Each translate is stable by Corollary \ref{cor:phase-monotone-stability}. Smooth convergence passes stability to $u_\pm$, regarded as functions on $\mathbb R^n$. To obtain stability on $\mathbb R^{n-1}$, test with $\phi(x')\eta(x_n/R)$, divide by the one-dimensional cutoff's squared integral, and let $R\to\infty$. Its gradient-to-mass ratio tends to zero like $R^{-2}$, leaving exactly the lower-dimensional quadratic form. This proves the end-limit assertion in \cite{ChAC}, Exercise 5.8, without assuming the ends are already the two wells.
\end{proof}

\begin{proposition}[Energy comparison with a monotone end limit]\label{prop:phase-translation-energy}
For $u$ as in Proposition \ref{prop:phase-monotone-end-limits} and any ball $B_R(a)\subset\mathbb R^n$,
\[
 E_1(u;B_R(a))\leq E_1(u_+;B_R(a))+C(n)R^{n-1},
\]
where $u_+(x')$ is regarded as a function on $\mathbb R^n$ independent of $x_n$.
\end{proposition}
\begin{proof}
Write $u_t(x)=u(x+te_n)$. Differentiate the energy on the fixed ball and use the equation to cancel its interior first variation:
\[
 \frac d{dt}E_1(u_t;B_R(a))
 =\int_{\partial B_R(a)}\partial_tu_t\,\partial_\nu u_t.
\]
The uniform gradient estimate and $\partial_tu_t>0$ bound this below by $-C(n)\int_{\partial B_R(a)}\partial_tu_t$. Integrate from zero to $T$ and let $T\to\infty$, using smooth convergence on the closed ball. Since $0\leq u_+-u\leq2$, the resulting boundary integral is bounded by a constant times the sphere area. This proves the energy-transport calculation in \cite{ChAC}, Section 5.5, p.~22.
\end{proof}

\begin{theorem}[Monotone bounded solutions in three dimensions]\label{thm:phase-monotone-three-dimensional}
Let $u:\mathbb R^3\to[-1,1]$ solve $\Delta u=W'(u)$ smoothly and satisfy $\partial_{x_3}u>0$. Then
\[
 u(x)=\tanh\left(\frac{e\cdot x-b}{\sqrt2}\right)
 \quad\text{with }|e|=1,\ e\cdot e_3>0.
\]
No prescribed limit as $x_3\to\pm\infty$ is assumed.
\end{theorem}
\begin{proof}
Strict monotonicity and boundedness give $-1<u<1$ and stability. The upper end limit is a bounded stable solution on $\mathbb R^2$ by Proposition \ref{prop:phase-monotone-end-limits}. Corollary \ref{cor:phase-planar-stable} makes it a constant well or a planar heteroclinic. Regarded on $\mathbb R^3$, either has energy in $B_R$ at most $c_W\pi R^2$: for a layer, slice in its transition direction, bound each transverse disk area by $\pi R^2$, and integrate its one-dimensional density of total mass $c_W$.

Proposition \ref{prop:phase-translation-energy} now gives $E_1(u;B_R)\leq CR^2$. Theorem \ref{thm:phase-stable-growth-rigidity} applies. The constant alternatives are excluded by strict monotonicity, which also forces the displayed direction to have positive third component. This is the argument of \cite{ChAC}, Theorem 5.3 and Section 5.5, with the end-limit stability and growth steps separated.
\end{proof}
